\documentclass[11pt]{exam}
\usepackage{foerderkurs}

\begin{document}

\minilektion{0.1}{Das Gleichheitszeichen}
  {das Gleichheitszeichen korrekt verwenden. Es steht nur zwischen zwei
   Ausdrücken, die denselben Wert haben.}

\begin{lehrperson}[Ablauf]
\textbf{Dauer:} ca.\ 12 Minuten Instruktion, danach Übungen in der Übungszeit,
Exit-Frage am Ende der Lektion.

\begin{tabular}{@{}ll@{}}
  2 min & Einstieg: Kästchenaufgabe verdeckt lösen lassen \\
  3 min & Antworten sammeln und besprechen, dann die falsche Kette zeigen \\
  4 min & Merkkasten und Beispiele 1–3 \\
  3 min & Andere Schreibweisen, Auftrag für die Übungen
\end{tabular}

\textbf{Material:} Mini-Whiteboards oder Zettel für die verdeckten Antworten.

\textbf{Verknüpfung:} Standortbestimmung 6.1, 6.2, 6.5\,b; Übungen mit den
Tags 0A und 6A.
\end{lehrperson}

% ---------------------------------------------------------------------------

\begin{einstieg}
\begin{enumerate}
  \item Welche Zahl gehört ins Kästchen? Schreiben Sie Ihre Antwort auf, ohne
        sie jemandem zu zeigen.
        \[ 8 + 4 = \kaestchen + 5 \]
  \item Lea hat die Aufgabe «Berechnen Sie $(3 + 4) \cdot 2 - 1$» so gelöst:
        \[ 3 + 4 = 7 \cdot 2 = 14 - 1 = 13 \]
        Ist das richtig aufgeschrieben? Ist $3 + 4$ gleich $13$?
\end{enumerate}
\end{einstieg}

\begin{lehrperson}[Einstieg]
Typische Antworten bei 1: \textbf{12}, \textbf{17} und \textbf{7}.
\begin{itemize}
  \item \textbf{12} ist die wichtigste falsche Antwort: Das Gleichheitszeichen
        wird als «jetzt kommt das Ergebnis» gelesen.
  \item \textbf{17} entsteht, wenn alle Zahlen zusammengezählt werden.
  \item Nicht sofort auflösen, sondern fragen: «Wie haben Sie gedacht?» Erst
        danach die Bedeutung «hat denselben Wert wie» einführen.
\end{itemize}
Bei 2: Die meisten Schülerinnen und Schüler wissen, was Lea meint. Der Punkt
ist: Lea hat richtig \emph{gedacht}, aber falsch \emph{geschrieben}. Mathematik
ist eine Sprache, und wer die Rechnung liest, liest wörtlich
«$3 + 4 = 13$». Das gilt auch für Lehrpersonen an Prüfungen.
\end{lehrperson}

% ---------------------------------------------------------------------------

\begin{merke}
\begin{itemize}
  \item Das Gleichheitszeichen bedeutet \textbf{«hat denselben Wert wie»},
        nicht «jetzt kommt das Ergebnis».
  \item Eine Gleichung gilt in \textbf{beide Richtungen}:
        $12 = 5x + 2$ bedeutet dasselbe wie $5x + 2 = 12$.
  \item In einer \textbf{Gleichungskette} $A = B = C$ müssen \textbf{alle}
        Teile denselben Wert haben.
  \item Eine \textbf{neue Rechnung} beginnt auf einer \textbf{neuen Zeile}.
  \item Beim Lösen von Gleichungen steht \textbf{pro Zeile eine Gleichung}.
        Die Umformung wird rechts mit einem senkrechten Strich angegeben,
        z.\,B.\ $\big|\; -2$.
\end{itemize}
\end{merke}

% ---------------------------------------------------------------------------
\abschnitt{Beispiele}

\begin{beispiel}[1: Eine korrekte Gleichungskette]
Berechnen Sie $3 \cdot (4 + 2)$.
\[ 3 \cdot (4 + 2) = 3 \cdot 6 = 18 \]
Jeder Teil der Kette hat den Wert $18$. Die Kette ist korrekt.
\end{beispiel}

\begin{beispiel}[2: Eine falsche Kette verbessern]
Berechnen Sie $(4 + 2) \cdot 3 - 5$.

\falsch\quad $4 + 2 = 6 \cdot 3 = 18 - 5 = 13$
\hfill (das würde heissen: $4 + 2 = 13$)

\richtig\quad als Kette, die immer mit dem ganzen Ausdruck weiterrechnet:
\[ (4 + 2) \cdot 3 - 5 = 6 \cdot 3 - 5 = 18 - 5 = 13 \]
\richtig\quad oder in getrennten Zeilen:
\[ 4 + 2 = 6 \qquad\qquad 6 \cdot 3 = 18 \qquad\qquad 18 - 5 = 13 \]
\end{beispiel}

\begin{beispiel}[3: Eine Gleichung lösen]
Lösen Sie $5x + 2 = 12$.

\falsch\quad $5x + 2 = 12 - 2 = 10 \geteilt 5 = 2$

\richtig\quad pro Zeile eine Gleichung, die Umformung steht rechts:
\begin{align*}
  5x + 2 &= 12 \umf{-2} \\
  5x     &= 10 \umf{\geteilt 5} \\
  x      &= 2
\end{align*}
Die Gleichheitszeichen stehen untereinander. So sieht man sofort, was
von Zeile zu Zeile geschieht.
\end{beispiel}

\begin{beispiel}[4: Gleichheit erkennen, ohne zu rechnen]
Stimmt $37 + 25 = 36 + 26$?

Ja: Links ist der erste Summand um $1$ grösser, der zweite um $1$ kleiner. Die
Summe bleibt gleich. Wer das Gleichheitszeichen als Beziehung versteht, muss
hier gar nicht rechnen.
\end{beispiel}

\begin{lehrperson}[Beispiele]
Beispiel 2 ist der Kern der Lektion. Es lohnt sich, die falsche Kette laut
vorzulesen: «Vier plus zwei ist gleich dreizehn.» Viele merken erst dabei,
was sie eigentlich aufgeschrieben haben.

Beispiel 3 nimmt die Darstellung beim Gleichungslösen vorweg (Mini-Lektion
6.1). Hier geht es nur um die Form, nicht um das Verfahren. Falls an Ihrer
Schule eine andere Schreibweise für die Umformung üblich ist (z.\,B.\ mit
$\Leftrightarrow$), diese als Alternative zeigen.

Beispiel 4 bereitet Aufgaben zum relationalen Denken vor (Übung 2).
\end{lehrperson}

% ---------------------------------------------------------------------------

\begin{schreibweisen}
\begin{itemize}
  \item $\approx$ bedeutet «ungefähr gleich»: $\pi \approx 3.14$.
  \item $\widehat{=}$ bedeutet «entspricht»: Auf einer Karte gilt
        $1\text{ cm} \mathrel{\widehat{=}} 50\text{ m}$. Das ist \emph{keine}
        Gleichheit, denn $1$\,cm ist nicht gleich $50$\,m.
  \item $\Leftrightarrow$ («ist gleichwertig mit») steht in manchen Büchern und
        Videos zwischen den Zeilen einer Gleichungsumformung.
  \item In Programmiersprachen wie Python bedeutet \texttt{x = x + 1} eine
        \emph{Zuweisung}: «Nimm den Wert von $x$, zähle $1$ dazu und speichere
        das Ergebnis wieder in $x$.» Als Gleichung wäre $x = x + 1$ unlösbar.
        Der Vergleich «ist gleich?» wird dort \texttt{==} geschrieben.
  \item In Antworten von KI-Programmen stehen zwischen Rechenschritten oft
        Pfeile wie \texttt{=>} oder \texttt{->}. Sie bedeuten «daraus folgt»
        und sind kein Gleichheitszeichen.
\end{itemize}
\end{schreibweisen}

% ---------------------------------------------------------------------------
\abschnitt{Übungen}

\begin{questions}

\question Welche Zahl gehört ins Kästchen?
\begin{multicols}{2}
\begin{parts}
  \part $15 + 7 = \kaestchen + 9$
  \part $6 \cdot 4 = 3 \cdot \kaestchen$
  \part $50 - 18 = \kaestchen - 20$
  \part $\kaestchen \geteilt 4 = 36 \geteilt 12$
\end{parts}
\end{multicols}
\begin{solution}
  a) $13$ \quad b) $8$ \quad c) $52$ \quad d) $12$
\end{solution}
\raum{2cm}

\question Ist die Aussage wahr oder falsch? Entscheiden Sie \textbf{ohne}
die Ergebnisse auszurechnen, und begründen Sie kurz.
\begin{parts}
  \part $48 + 27 = 50 + 25$     \hfill \wahlr{wahr}\quad \wahl{falsch}
  \part $63 - 29 = 64 - 30$     \hfill \wahlr{wahr}\quad \wahl{falsch}
  \part $8 \cdot 25 = 4 \cdot 50$ \hfill \wahlr{wahr}\quad \wahl{falsch}
  \part $72 \geteilt 6 = 36 \geteilt 12$ \hfill \wahl{wahr}\quad \wahlr{falsch}
  \part $99 + 46 = 100 + 47$    \hfill \wahl{wahr}\quad \wahlr{falsch}
\end{parts}
\begin{solution}
  a) wahr: $+2$ und $-2$ gleichen sich aus. \
  b) wahr: Bei der Subtraktion müssen \emph{beide} Zahlen um gleich viel in dieselbe Richtung
  verändert werden. \
  c) wahr: ein Faktor halbiert, der andere verdoppelt. \
  d) falsch: Halbiert man den Dividenden, muss man auch den Divisor
  halbieren ($36 \geteilt 3$), nicht verdoppeln. \
  e) falsch: Beide Summanden werden grösser.
\end{solution}
\raum{3cm}

\question Die folgenden Rechnungen haben das richtige Ergebnis, sind aber
falsch aufgeschrieben. Finden Sie den Fehler und schreiben Sie die Rechnung
korrekt auf.
\begin{parts}
  \part Aufgabe: $(12 \geteilt 3 + 5) \cdot 2$. \quad Lösung:
        $12 \geteilt 3 = 4 + 5 = 9 \cdot 2 = 18$
  \part Aufgabe: $(2 + 3)^2$. \quad Lösung: $2 + 3 = 5^2 = 25$
  \part Aufgabe: Lösen Sie $3x - 4 = 11$. \quad Lösung:
        $3x - 4 = 11 + 4 = 15 \geteilt 3 = 5$
\end{parts}
\begin{solution}
  a) $(12 \geteilt 3 + 5) \cdot 2 = (4 + 5) \cdot 2 = 9 \cdot 2 = 18$ \quad
  b) $(2 + 3)^2 = 5^2 = 25$ \quad
  c) $3x - 4 = 11 \;\big|\; +4$, \ $3x = 15 \;\big|\; \geteilt 3$, \ $x = 5$
  (untereinander geschrieben)
\end{solution}
\raum{5cm}

\question Lösen Sie die Gleichungen. Schreiben Sie pro Zeile eine Gleichung
und geben Sie die Umformungen rechts an.
\begin{multicols}{3}
\begin{parts}
  \part $4x + 7 = 19$
  \part $20 = 3x - 1$
  \part $2x - 5 = 9$
\end{parts}
\end{multicols}
\begin{solution}
  a) $x = 3$ \quad b) $x = 7$ \quad c) $x = 7$. \ Bei b) darf die Gleichung
  zuerst umgedreht werden ($3x - 1 = 20$), muss aber nicht.
\end{solution}
\raum{5cm}

\end{questions}

\begin{lehrperson}[Übungen]
Übung 3 ist die wichtigste: Die Ergebnisse stimmen, nur die Darstellung ist
falsch. Das zeigt, dass es hier nicht ums Rechnen geht. Übung 4 erst nach
Übung 3 lösen lassen. Wer bei Übung 2 d) «wahr» ankreuzt, überträgt die
Regel der Multiplikation unbesehen auf die Division; das in Mini-Lektion 2.1
wieder aufgreifen.
\end{lehrperson}

% ---------------------------------------------------------------------------

\begin{ausblick}
Bei jeder Herleitung, bei jedem Gleichungssystem und bei jedem Beweis werden
Sie Gleichungen und Gleichungsketten aufschreiben. Eine falsche Kette gibt an
Prüfungen Abzug, auch wenn das Ergebnis stimmt. Und in der Informatik
begegnet Ihnen das Gleichheitszeichen mit einer ganz anderen Bedeutung wieder.
\end{ausblick}

\begin{exitfrage}
Was ist an «$5 \cdot 4 = 20 + 3 = 23$» falsch? Schreiben Sie die Rechnung
korrekt auf.
\ifprintanswers\par\medskip\textbf{Lösung:} Es steht $5 \cdot 4 = 23$.
Korrekt: $5 \cdot 4 + 3 = 20 + 3 = 23$.\else\par\vspace{2cm}\fi
\end{exitfrage}

\end{document}
